%!TEX root=../mythesis.tex
% Chapter 2 -- Preliminaries

\chapter{Preliminaries} % Main chapter title
\chaptermark{Preliminaries}
\label{ch:preliminaries}

This chapter introduces the shared vocabulary used throughout the thesis. \sref{sec:bg_finance} introduces the financial instruments, market microstructure, and evaluation metrics common to all three studies; the metric definitions given there are referenced rather than repeated in later chapters. \sref{sec:bg_rl} is a brief primer on the reinforcement-learning machinery---from Markov decision processes to GFlowNets and quality-diversity search---used across the thesis. The survey of related work and the open challenges that motivate the technical chapters are deferred to \cref{ch:literature_review}.

%----------------------------------------------------------------------------------------
\section{Financial Markets and Quantitative Trading}
\label{sec:bg_finance}

\subsection{Instruments and Market Microstructure}
\label{sec:bg_instruments}

The three studies operate on three increasingly complex trading settings. EarnHFT (\cref{ch:earnhft}) trades a single \emph{spot} cryptocurrency, where the asset is bought and sold outright and only long positions are taken. FineFT (\cref{ch:fineft}) trades \emph{perpetual futures}, standardised contracts that permit both long and short positions, magnify exposure through leverage, and exchange a periodic funding payment between the two sides of the contract; depending on the sign of the funding rate and the direction of the position, a trader may either pay or receive funding. AlphaQD (\cref{ch:alphaqd}) operates on the \emph{cross-section} of an entire equity market, where the object of interest is the relative ranking of thousands of assets on each day rather than the timing of a single instrument. The following concepts recur throughout.

\paragraph{Limit order book (LOB).} The limit order book aggregates all resting buy and sell orders by price level and is the standard description of market microstructure~\cite{madhavan2000market}. A $d$-level LOB at time $t$ is denoted $B_t=(p_t^{b_1},p_t^{a_1},q_t^{b_1},q_t^{a_1},\dots,p_t^{b_d},p_t^{a_d},q_t^{b_d},q_t^{a_d})$, where $p_t^{b_i}$ and $p_t^{a_i}$ are the level-$i$ bid and ask prices and $q_t^{b_i},q_t^{a_i}$ are the corresponding quantities (the symbol $m$ is reserved for the annualisation constant of \sref{sec:bg_metrics}). \fref{fig:bg_lob} illustrates a snapshot. A \emph{market order} executes immediately against the resting orders on the opposite side, filling at the best quote first; only when its size exceeds the depth available at the best level does it \emph{walk the book}, consuming successively deeper levels, so that its average execution price is worse than the best quote---an effect known as slippage that any high-fidelity simulator must model.

\begin{figure}[!htb]
  \centering
  \begin{tabular}{lrr}
    \toprule
    Level & Price & Quantity \\
    \midrule
    Ask 5 & 105 & 90 \\
    Ask 4 & 104 & 62 \\
    Ask 3 & 103 & 45 \\
    Ask 2 & 102 & 28 \\
    Ask 1 (best ask) & 101 & 10 \\
    \midrule
    \multicolumn{3}{c}{mid-price $=100.5$} \\
    \midrule
    Bid 1 (best bid) & 100 & 12 \\
    Bid 2 & 99 & 30 \\
    Bid 3 & 98 & 48 \\
    Bid 4 & 97 & 66 \\
    Bid 5 & 96 & 85 \\
    \bottomrule
  \end{tabular}
  \caption[A limit order book snapshot.]{A schematic five-level limit order book snapshot with illustrative integer prices and quantities. Resting asks (sell orders) sit above the mid-price and resting bids (buy orders) below it, each aggregated by price level. A market order first consumes the depth at the best quote; if its size exceeds that depth, it walks the book \emph{away} from the mid-price into successively deeper levels, so larger orders execute at progressively worse average prices.}
  \label{fig:bg_lob}
\end{figure}

\paragraph{OHLC and trades.} Executed transactions are summarised over a time interval by the open, high, low, and close prices, denoted $x_t=(p_t^o,p_t^h,p_t^l,p_t^c)$. The raw stream of executed orders, or trades, records price, quantity, and side. For leveraged contracts, a \emph{mark price} $M_t$ is used as the reference fair value at which the position is valued and against which liquidation is assessed.

\paragraph{Technical indicators.} Both microstructure and price-history features are summarised by technical indicators---deterministic functions of recent LOB and trade data designed to expose latent patterns, denoted $y_t=\phi(\cdot)$. Examples are the imbalance volume (IV)~\cite{chordia2002order}, which measures buying versus selling pressure in the book, and the moving-average convergence divergence (MACD)~\cite{appel2005technical}, which contrasts trends over different windows. These indicators are inputs to the RL state in the two trading studies (\cref{ch:earnhft} and \cref{ch:fineft}); the alpha-mining study of \cref{ch:alphaqd} instead composes its own signals directly from raw price--volume fields. As standalone rules, however, technical indicators are sensitive to hyperparameters and produce false signals in volatile markets.

\paragraph{Position, leverage, and value.} A trader's position $H_t$ is the signed quantity held; under leverage $L_t$, the same capital supports a larger position and amplifies both gains and losses. For spot trading, the account value (net value) is the cash balance plus the position marked at the current price, net of realised costs such as commissions. For leveraged futures, the corresponding quantity is the \emph{margin balance}: the wallet balance plus the unrealised profit-and-loss of open positions, net of commissions and funding payments---the full notional value of the position does not enter the account value, which is precisely what leverage means.

\subsection{The Trading-Frequency Spectrum}
\label{sec:bg_frequency}

Trading strategies are conventionally classified by the frequency at which they act, and the binding difficulty changes along this spectrum. \emph{High-frequency} strategies act at the second or sub-second scale and profit from fleeting microstructure imbalances; their defining computational challenge is the sheer length of the decision horizon. \emph{Medium-frequency} strategies act at the minute-to-hour scale. \emph{Low-frequency} strategies act daily or slower; here the difficulty is not the horizon but the discovery of robust predictive signals from a low signal-to-noise cross-section. Acting frequency should not be conflated with other, logically independent dimensions of a trading problem: leverage, funding payments, and overnight exposure are properties of the instrument traded and of the holding period, not of the sampling rate itself, although longer holding periods tend to make them more consequential. In this thesis these dimensions become material in the leveraged perpetual-futures setting of \cref{ch:fineft}. Note also that the 24/7 operation of cryptocurrency markets removes the closed-market gap between trading sessions, but not the risk of holding a position overnight. This thesis traverses the entire frequency spectrum, from second-level execution in \cref{ch:earnhft} to daily signal generation in \cref{ch:alphaqd}.

\subsection{Evaluation Metrics}
\label{sec:bg_metrics}

Strategies are evaluated at two layers. The \emph{portfolio layer} measures the realised profit-and-loss of an executed strategy; the \emph{signal layer}, used for alpha mining, measures the predictive quality of a signal independently of any portfolio construction. Let $V_1,\dots,V_{n+1}$ denote the account value at the $n{+}1$ decision points of the trading period, and let $\mathbf{ret}=(ret_1,\dots,ret_n)$ with $ret_i=(V_{i+1}-V_i)/V_i$ denote the per-step \emph{simple} (not log) returns. Throughout, $m$ is the number of decision steps in a year at the strategy's acting frequency (the annualisation constant), the risk-free rate is taken to be zero, and $\mathbb{E}[\cdot]$ and $\sigma[\cdot]$ denote the sample mean and sample standard deviation over the return series.

\paragraph{Portfolio-layer metrics.} The following six metrics are used throughout \cref{ch:earnhft} and \cref{ch:fineft}, and the return/risk pair recurs in \cref{ch:alphaqd}.
\begin{itemize}
\item \textbf{Total Return (TR)}: the overall return over the trading period, $TR=(V_{n+1}-V_1)/V_1$, where $V_1$ and $V_{n+1}$ are the initial and final account values.
\item \textbf{Annualised Volatility (AVOL)}: $\mathrm{AVOL}=\sigma[\mathbf{ret}]\,\sqrt{m}$, a measure of risk.
\item \textbf{Maximum Drawdown (MDD)}~\cite{magdon2004maximum}: the largest peak-to-trough decline of the value curve, $\mathrm{MDD}=\max_{1\le i\le j\le n+1}\,(V_i-V_j)/V_i$, capturing the worst case. MDD is reported as a positive fraction throughout this thesis, so smaller is better (marked $\downarrow$ in the result tables).
\item \textbf{Annualised Sharpe Ratio (ASR)}~\cite{sharpe1994sharpe}: $\mathrm{ASR}=\sqrt{m}\,\frac{\mathbb{E}[\mathbf{ret}]}{\sigma[\mathbf{ret}]}$, return per unit of total risk (with zero risk-free rate); the $\sqrt{m}$ factor multiplies the whole ratio, so ASR equals the annualised mean return $m\,\mathbb{E}[\mathbf{ret}]$ divided by AVOL.
\item \textbf{Annualised Calmar Ratio (ACR)}: $\mathrm{ACR}=m\,\mathbb{E}[\mathbf{ret}]/\mathrm{MDD}$, return per unit of worst-case loss.
\item \textbf{Annualised Sortino Ratio (ASoR)}: $\mathrm{ASoR}=\sqrt{m}\,\frac{\mathbb{E}[\mathbf{ret}]}{DD}$, where the downside deviation $DD$ is the standard deviation of the negative per-step returns (downside threshold zero) and the $\sqrt{m}$ factor is placed as in the ASR; it penalises only harmful volatility.
\end{itemize}

\paragraph{Signal-layer metrics.} For a cross-sectional signal $\alpha$ and an $h$-day forward return $r^{(h)}$, the daily \emph{information coefficient} is the cross-sectional correlation $\mathrm{IC}_t(\alpha)=\corr(\alpha_t,r^{(h)}_t)$. Its time average and stability are combined into the \emph{information ratio} $\mathrm{ICIR}(\alpha)=\overline{\mathrm{IC}(\alpha)}/\hat\sigma(\mathrm{IC}(\alpha))$, the canonical scalar quality of a single alpha in the active-management literature~\cite{grinold2000active} and standard engineering practice in formulaic-alpha libraries~\cite{kakushadze2016101}, since it rewards predictive strength and temporal stability simultaneously. The rank-based variant, \emph{RankIC}, replaces Pearson with Spearman correlation and is robust to outliers.

%----------------------------------------------------------------------------------------
\section{Reinforcement Learning Foundations}
\label{sec:bg_rl}

This section is a brief primer, confined to the reinforcement-learning machinery the technical chapters actually use; a reader fluent in RL may skip directly to the task survey of \cref{ch:literature_review}, which carries the weight of the literature review.

\subsection{Markov Decision Processes}
\label{sec:bg_mdp}

A Markov decision process (MDP) is the tuple $(\mathcal{S},\mathcal{A},P,r,\gamma,T)$, where $\mathcal{S}$ is the state space, $\mathcal{A}$ the action space, $P:\mathcal{S}\times\mathcal{A}\times\mathcal{S}\to[0,1]$ the transition function, $r$ the reward function, $\gamma\in(0,1]$ the discount factor, and $T$ the horizon. At each step the agent observes $s_t$, takes an action $a_t$ under a policy $\pi_\theta:\mathcal{S}\times\mathcal{A}\to[0,1]$, and receives a reward $r_t$ and a next state $s_{t+1}$. The objective is to find parameters maximising the expected discounted return, $\theta^\star=\argmax_\theta \mathbb{E}_{s_0\sim\rho_0,\,\pi_\theta}\big[\sum_{t=0}^{T}\gamma^t r_t\big]$, where $\rho_0$ is the initial-state distribution; the optimal policy is then $\pi^\star=\pi_{\theta^\star}$. The convergence guarantees of standard RL algorithms presuppose that $P$ and $r$ are \emph{stationary}; the central difficulty of this thesis is that financial environments are not, a fact made explicit in \cref{ch:fineft} by the dynamic-parametric MDP formulation that augments the tuple with a slowly evolving latent dynamic.

\subsection{Value-Based Methods}
\label{sec:bg_value}

Value-based methods learn action-value functions. For a policy $\pi$, $Q^\pi(s,a)$ is the expected return of taking $a$ in $s$ and following $\pi$ thereafter; the \emph{optimal} action-value function $Q^\star(s,a)=\max_\pi Q^\pi(s,a)$ is the expected return of taking $a$ in $s$ and acting optimally thereafter. Q-learning~\cite{watkins1992qlearning} learns an estimate $Q$ of $Q^\star$ by updating it towards the temporal-difference (TD) target $r+\gamma\max_{a'}Q(s',a')$. The Deep Q-Network (DQN)~\cite{mnih2015human} parameterises $Q$ with a neural network and stabilises learning through experience replay and a target network. Double DQN (DDQN)~\cite{van2016deep} decouples action selection from evaluation to reduce the overestimation bias of the $\max$ operator. EarnHFT and FineFT both build on DDQN, augmenting its TD target with supervision from a dynamic-programming teacher, a design related in spirit to deep Q-learning from demonstrations~\cite{hester2018dqfd}; the convergence of this augmented operator is established in \cref{ch:earnhft}.

\subsection{Policy-Based Methods}
\label{sec:bg_policy}

Policy-based methods optimise the policy directly. Proximal Policy Optimisation (PPO)~\cite{schulman2017proximal} maximises a clipped surrogate objective with importance sampling, trading some bias for stability and sample reuse. In trading, policy-based methods are prone to collapsing onto trivial policies---holding the current position---when transaction costs make any action locally unprofitable, a failure mode documented empirically in \cref{ch:earnhft}.

\subsection{Distributional and Ensemble Reinforcement Learning}
\label{sec:bg_dist_ensemble}

Two families address the high stochasticity of financial rewards. \emph{Distributional} RL learns the full return distribution rather than its mean; quantile-regression DQN~\cite{dabney2018distributional} and its risk-averse variants represent the return by a set of quantiles and can act on the lower tail to avoid worst cases. \emph{Ensemble} RL trains several learners and aggregates them: averaged DQN reduces target variance, while SUNRISE~\cite{lee2021sunrise} treats inter-learner disagreement as a measure of uncertainty and down-weights uncertain samples. FineFT belongs to the ensemble family but departs from prior work by updating learners \emph{selectively} rather than uniformly, which is what allows specialists to emerge.

\subsection{Hierarchical Reinforcement Learning}
\label{sec:bg_hrl}

Hierarchical RL decomposes a long-horizon task into a hierarchy of sub-tasks, typically a high-level policy that selects among low-level policies or temporally extended options~\cite{sutton1999between,pateria2021hierarchical}. In trading, this naturally maps onto a separation of time scales: a slow policy that reacts to macro conditions and fast policies that execute. Prior hierarchical trading frameworks address portfolio management and order execution~\cite{wang2021commission} or select among strategies by market situation~\cite{niu2022metatrader}. EarnHFT brings this idea to single-asset high-frequency trading, where the high-level router operates at the minute scale and the low-level specialists at the second scale.

\subsection{Generative Flow Networks}
\label{sec:bg_gflownet}

Generative flow networks (GFlowNets)~\cite{bengio2021gflownet} learn a stochastic policy that constructs compositional objects---here, symbolic expressions---step by step, with the aim that the probability of producing a complete object be proportional to a terminal reward. This proportionality holds exactly only at the global optimum of the training objective; at the optimum of the trajectory-balance objective~\cite{malkin2022trajectorybalance} used here, the sampler is reward-proportional, and in practice a trained GFlowNet samples approximately proportionally to reward rather than collapsing onto the single maximiser, so it covers many high-reward modes of the search space under a fixed budget. This mode-covering property is precisely what an alpha-mining search requires, and AlphaQD adopts a relational-graph GFlowNet as its generator.

\subsection{Quality-Diversity Search and MAP-Elites}
\label{sec:bg_qd}

Quality-diversity (QD) algorithms~\cite{pugh2016quality} seek not a single optimum but an archive of high-quality solutions that are diverse along a chosen \emph{behaviour descriptor}. MAP-Elites~\cite{mouret2015mapelites} discretises the descriptor space into cells and retains the best solution(s) per cell, yielding a structured population that covers the behaviour space. QD has not previously been applied to formulaic alpha mining; AlphaQD's contribution is to identify a meaningful behaviour descriptor for this domain---the normalised daily IC series, an empirical proxy for a formula's return source---and to use the archive as the mechanism that organises diversity and selection outside the learning reward.

\subsection{Variational Autoencoders and Out-of-Distribution Detection}
\label{sec:bg_vae}

A variational autoencoder (VAE)~\cite{kingma2013auto} learns a probabilistic encoder that maps inputs to a regularised latent distribution and a decoder that reconstructs them. Inputs resembling the training distribution tend to reconstruct well, so the VAE reconstruction loss is widely used as a heuristic reconstruction-based score for out-of-distribution (OOD) detection; it is a practical signal rather than a calibrated probabilistic measure. FineFT trains a VAE per market dynamic and uses its loss to detect when the current market state lies outside a specialist's competence, converting OOD detection into an explicit risk control.
